\chapter{\(\mathrm{L}_{\mathrm{ACD}}\) Exceptional-Crossing \(\mathrm{Ext}^{1}\)-Class Match at \(\Sigma^{\mathrm{exc}}=\{0,1\}\)}
\label{v4-arith:w10-d:chapter}

\emph{The Langlands--Shahidi normalised intertwining operator
\(N(s,\chi)\) of \Cref{v4-arith:w6-d:def:N-normalised} has
a first-order pole at \(s=1\), while its functional-equation dual has
the corresponding first-order pole at \(s=0\). Chapter~\ref{v4-arith:w6-d:chapter}
reduced the \(\mathrm{L}_{\mathrm{ACD}}\) comparison on the
non-exceptional sph-finite locus
\(\Sigma^{\mathrm{nonexc}}=i\mathbb{R}\setminus\{0\}\) and stated
the two remaining points as Conjecture~\ref{v4-arith:w6-d:conj:exceptional-unitary-crossing}.
This chapter treats the exceptional class in that comparison under the
Emerton--Gee wild-cyclotomic enhancement
(Conjecture~\ref{v4-arith:deg-f1a:conj:EG-wild-enhancement}, the
same hypothesis sustaining the F1.a.2 class-match of
\Cref{v4-arith:w6-f:thm:class-match} at \(p\in\{2,3\}\)), and on the
restricted-product Paley--Wiener Ext comparison for the automorphic
\(\mathrm{Ext}^{1}\) side.
The argument proceeds through four modules. Module~(X0) records
the two-point residue data: at \(s=1\) the residual spectrum is the
one-dimensional trivial representation (Mo{\oe}glin--Waldspurger
II.2.4); at \(s=0\) the functional-equation dual is a reducible
spherical induction with Steinberg subquotient (Shahidi 2010 \S 3.5).
Module~(X1) computes the local \(\mathrm{Ext}^{1}\)-extension class
at each crossing from the first-order residue of \(N_{v}(s,\chi)\)
(Shahidi 2010 Thm~3.4.1) and assembles the adelic class via Tate
factorisation. Module~(X2) identifies the spectral
\(\mathrm{Ext}^{1}\)-class on the Emerton--Gee stack via Arthur 2013
\S 1 endoscopic classification: the residual summand at \(s=1\) lies
in the trivial Arthur packet and at \(s=0\) in the Steinberg Arthur
	packet; in both cases the Arthur \(\mathrm{SL}_{2}\)-factor determines
	a deformation-theoretic tangent direction matching the
\(\mathrm{Ext}^{1}\) of (X1). Module~(X3) establishes the
class-level match via the arithmetic Koszul--Petersson compatibility
of Chapter~\ref{v4-arith:w6-c:LS-correction}: the LS-corrected Koszul
pairing transports the appropriate leading LS coefficient onto the
	tangent direction at the Arthur stratum. The
Bushnell--Henniart block of Chapter~\ref{v4-arith:f1a3-BH} handles
the \(p\)-adic component; the archimedean component uses
Harish-Chandra 1976 Thm~1. Under these two hypotheses, the exceptional
\(\mathrm{Ext}^{1}\)-class is matched; the
global \(\mathrm{L}_{\mathrm{ACD}}\) comparison still depends on the
standing local-global and Bernstein-centre gluing inputs.}

\section{Objects and Sources}
\label{v4-arith:w10-d:domain}

\subsection{Inherited Objects}

The following objects from earlier chapters are used:
\begin{itemize}
\item the Langlands--Shahidi normalised intertwining operator
\(N_{v}(s,\chi)\) with explicit exceptional pole/zero divisor at \(s=1,0\)
(\Cref{v4-arith:w6-d:def:N-normalised},
\Cref{v4-arith:w6-d:prop:local-FE});
\item the exceptional locus
\(\Sigma^{\mathrm{exc}}=\{s=0\}\cup\{s=1\}\subset\mathbb{C}\)
(\Cref{v4-arith:w6-d:def:exceptional}) and its named class-match
conjecture (\Cref{v4-arith:w6-d:conj:exceptional-unitary-crossing});
\item the LS-corrected arithmetic Koszul pairing
\(\langle\,\cdot\,,\,\cdot\,\rangle^{\mathrm{Koszul},\mathrm{LS}}\)
(\Cref{v4-arith:w6-c:def:LS-Koszul}) and the unconditional density
match (\Cref{v4-arith:w6-c:cor:unconditional});
\item the F1.a.2 class-match pattern
(\Cref{v4-arith:w6-f:thm:class-match}) and the residual gerbe
tangent identification
(\Cref{v4-arith:w6-f:prop:gerbe-cotangent-match});
\item the Langlands spectral decomposition
\(L^{2}([GL_{2}])=L^{2}_{\mathrm{cusp}}\oplus L^{2}_{\mathrm{res}}\oplus L^{2}_{\mathrm{cont}}\)
(\Cref{v4-arith:w6-d:thm:Langlands-decomp}) and the categorical
lift of the residual summand
(\Cref{v4-arith:w6-d:thm:res-lift}).
\end{itemize}

No new categorical object is constructed here. The match at the
two crossings is determined by these inputs and by the primary-source
residue computations of Mo{\oe}glin--Waldspurger, Shahidi, and
Arthur.

\subsection{Primary Sources}
\label{v4-arith:w10-d:sources}

The primary sources are:
\begin{description}
\item[Mo{\oe}glin--Waldspurger 1995] \emph{Spectral Decomposition
and Eisenstein Series}, Cambridge Tracts~113, Ch.~II.1--II.2 and
IV.1. The residual spectrum of \(GL_{2}/\mathbb{Q}\): II.2.4 identifies
\(L^{2}_{\mathrm{res}}([GL_{2}])\) as the one-dimensional trivial
representation (the constant function); IV.1.8 computes the local
and global residues of \(N(s,\chi)\) at \(s=1\).
\item[Shahidi 2010] \emph{Eisenstein Series and Automorphic
\(L\)-Functions}, AMS Colloquium~58, Thm~3.4.1 and \S 3.5. The
polar divisor of the normalised intertwining operator and its
functional-equation dual: at the first-order singular point the
leading coefficient is a non-zero intertwining operator between the two
Langlands subquotients of \(I_{v}(s_{0},\chi)\).
\item[Arthur 2013] \emph{The Endoscopic Classification of
Representations}, AMS Colloquium~61, Ch.~1 \S 1.4--1.5 and \S 3.5.
The Arthur parameter of a residual representation: at \(s=1\) the
trivial character sits in the trivial Arthur packet
\(\psi_{\mathrm{triv}}=[\mathrm{triv}\boxtimes\nu_{2}]\) with
\(\mathrm{SL}_{2}\)-factor \(\nu_{2}\) the two-dimensional irreducible
of \(\mathrm{SL}_{2}(\mathbb{C})\); at \(s=0\) the dual packet contains
the Steinberg representation.
\item[Emerton--Gee arXiv:1908.07185] The Emerton--Gee stack.
\S 6.2 unramified-locus smoothness; \S 6.3 the closed point at the
trivial character (cyclotomic or trivial); the formal completion at
the cyclotomic residual Galois representation is the universal
deformation ring of the reducible residual with cyclic-Weyl
structure.
\item[Breuil--Mezard 2002] Duke Math.~J.~115. Local deformation
multiplicities for two-dimensional Galois representations. The
one-dimensionality of the global residual spectrum used here is the
Mo{\oe}glin--Waldspurger statement above.
\item[Harish-Chandra 1976] Invent.~Math.~36. Archimedean Plancherel
density; the archimedean component of the residue at \(s_{0}=1\)
computed via \(\Gamma_{\mathbb{R}}(2s-1)/\Gamma_{\mathbb{R}}(2s)\)
has a simple pole with explicit residue.
\item[Schlessinger 1968, Mazur 1989] Deformation-theoretic
representability; the tangent space of the universal deformation
ring at a reducible residual equals
\(H^{1}(G_{\mathbb{Q}},\mathrm{ad}^{0}\bar\rho)\).
\item[Bella{\"i}che--Chenevier 2009] Ast{\'e}risque 324, Ch.~2.
Residual gerbe at a reducible closed point; the automorphism group
contributing to the gerbe tangent complex.
\end{description}

Secondary: Langlands 1976 \S 6--7 (original Eisenstein functional
equation and residue analysis); Kisin 2008 J.~AMS~21 (cotangent
complex of local Galois deformation rings); Tate 1950 Ch.~XV
(adelic restricted tensor product factorisation).

\subsection{Parallelism with F1.a.2.class-match}

The match pattern is:
\[
\text{automorphic }\mathrm{Ext}^{1}\;\xrightarrow{\;\text{connecting map}\;}\;\text{Galois }H^{1}\;\xrightarrow{\;\text{gerbe tangent}\;}\;\text{spectral }\mathrm{Ext}^{1}.
\]
We write
\[
\mathbb{T}^{1}_{\mathcal X}|_{x}:=
H^{0}\bigl((\mathbb L_{\mathcal X}|_{x})^{\vee}\bigr)
\]
for the tangent fibre used below.
F1.a.2 used the Pa{\v s}k{\=u}nas connecting map (non-split
block extensions over the Iwahori algebra) against the residual
gerbe tangent of the wild-enhanced Emerton--Gee stack at
\(p\in\{2,3\}\)
(\Cref{v4-arith:w6-f:prop:Paskunas-equals-Galois}--\ref{v4-arith:w6-f:prop:gerbe-cotangent-match}).
The exceptional crossing at \(s=1\) (respectively \(s=0\)) uses the
Mo{\oe}glin--Waldspurger + Shahidi residue at \(s=1\) (respectively
the functional-equation dual) against the residual gerbe tangent
at the cyclotomic character \(\chi_{\mathrm{cyc}}\) (respectively the
trivial character), with Arthur 2013 endoscopic packet as the
spectral-side intermediary. The two matches are formally parallel;
the residue at an \(N(s,\chi)\) pole plays the role of the Pa{\v
s}k{\=u}nas connecting map, and the Arthur packet at the
cyclotomic character plays the role of the BH type at the wild
prime.

\subsection{Independence of Sources}

Proof sources for Chapter~\ref{v4-arith:w6-d:chapter}
(Langlands 1976, Shahidi 1990/2010, Mo{\oe}glin--Waldspurger II/IV,
Harish-Chandra 1976, Wallach 1992, Arthur 2013 \S 3.5) deliver the
non-exceptional categorical equivalence and isolate the two-point
residual. Verification sources here are
Mo{\oe}glin--Waldspurger II.2 (residual spectrum, the constant
function), Shahidi 2010 Thm~3.4.1 (residue of \(N\) at a first-order
pole), Arthur 2013 \S 1.4--1.5 (Arthur packets for residual
representations), Emerton--Gee \S 6.2--6.3 (spectral stack at the
cyclotomic closed point), Breuil--Mezard 2002 (local deformation
multiplicities), Bella{\"i}che--Chenevier 2009 (residual
gerbe at a reducible closed point), Schlessinger 1968 + Mazur 1989
(deformation-theoretic representability and tangent space).
Mo{\oe}glin--Waldspurger, Shahidi 2010, and Arthur 2013 appear on
both sides of the disjointness partition, but in disjoint chapters:
Chapter~\ref{v4-arith:w6-d:chapter} used
Mo{\oe}glin--Waldspurger~IV (Eisenstein synthesis), Shahidi 2010 \S 3.5
(polar analysis), Arthur 2013 \S 3.5 (Paley--Wiener); this chapter
uses Mo{\oe}glin--Waldspurger~II.2 (residual spectrum), Shahidi
2010 Thm~3.4.1 (residue computation), Arthur 2013 \S 1.4--1.5
(endoscopic packets). The Emerton--Gee, Breuil--Mezard,
Bella{\"i}che--Chenevier, Schlessinger, and Mazur inputs are
spectral-side / deformation-theoretic and do not appear in the
Eisenstein / Paley--Wiener derivation set. Source independence
passes.

\section{Module (X0): Two-Point Residue Data on the Automorphic Side}
\label{v4-arith:w10-d:X0}

\subsection{The Residual Spectrum at \(s=1\)}

\begin{proposition}[residual spectrum of \(GL_{2}/\mathbb{Q}\) at \(s=1\)]
\label{v4-arith:w10-d:prop:residual-s1}
\ClaimStatusProvedElsewhere. The residual summand of
\(L^{2}([GL_{2}])\) at \(s=1\) is one-dimensional and spanned by the
constant function
\(\mathbb{1}_{[GL_{2}]}\colon[GL_{2}]\to\mathbb{C}\), \(g\mapsto 1\):
\begin{equation}
\label{v4-arith:w10-d:eq:residual-s1}
L^{2}_{\mathrm{res},s=1}([GL_{2}])\;=\;\mathbb{C}\cdot\mathbb{1}_{[GL_{2}]}.
\end{equation}
Equivalently, the spherical Eisenstein series
\(E(s,g)=\sum_{\gamma\in B(\mathbb{Q})\backslash GL_{2}(\mathbb{Q})}f_{s}(\gamma g)\)
(where \(f_{s}\) is the standard spherical section of \(I(s,\chi_{\mathrm{triv}})\))
has a simple pole at \(s=1\) with residue a constant function:
\begin{equation}
\label{v4-arith:w10-d:eq:residue-E}
\mathrm{Res}_{s=1}E(s,g)\;=\;c_{1}\cdot\mathbb{1}_{[GL_{2}]},\qquad c_{1}\in\mathbb{C}^{\times}.
\end{equation}
\end{proposition}

\begin{proof}[Source]
Mo{\oe}glin--Waldspurger 1995 Ch.~II.2.4 computes the residual
spectrum of \(GL_{n}/\mathbb{Q}\) and establishes the one-dimensional
trivial-character description for \(n=2\). The residue of the
Eisenstein series at \(s=1\) is the constant function
(Langlands 1976 \S 6; Mo{\oe}glin--Waldspurger II.2.4). The scalar
\(c_{1}\) is determined by the Maass--Selberg formula
(Mo{\oe}glin--Waldspurger II.2.4 and IV.1.8). Only its non-vanishing is
used below.
\end{proof}

\subsection{The Spherical Reducibility at \(s=0\)}

\begin{proposition}[reducibility at \(s=0\)]
\label{v4-arith:w10-d:prop:reducibility-s0}
\ClaimStatusProvedElsewhere. At \(s=0\), the spherical induction
\(I_{v}(0,\chi_{\mathrm{triv}})=\mathrm{Ind}_{B}^{GL_{2}}(1\otimes 1)\)
is reducible with exactly two composition factors:
\begin{equation}
\label{v4-arith:w10-d:eq:reducibility-s0}
0\;\longrightarrow\;\mathrm{St}_{v}\;\longrightarrow\;I_{v}(0,\chi_{\mathrm{triv}})\;\longrightarrow\;\mathbf{1}_{v}\;\longrightarrow\;0,
\end{equation}
where \(\mathrm{St}_{v}\) is the Steinberg representation
(one-dimensional twist of the special representation at
archimedean~\(v\); infinite-dimensional discrete-series at
non-archimedean~\(v\)) and \(\mathbf{1}_{v}\) is the trivial
representation. The short exact sequence defines a non-trivial
\(\mathrm{Ext}^{1}\)-class
\(\alpha_{v}^{s=0}\in\mathrm{Ext}^{1}(\mathbf{1}_{v},\mathrm{St}_{v})\).
\end{proposition}

\begin{proof}[Source]
Shahidi 2010 \S 3.5, Thm~3.5.1: the principal-series \(I_{v}(s,\chi)\)
for \(GL_{2}\) is reducible exactly at \(s=0,1\) with the listed
composition factors. The non-splitting of the sequence is
Bernstein--Zelevinsky 1977 \S 4 (non-archimedean) and
Jacquet--Langlands 1970 \S 5 (archimedean).
\end{proof}

\subsection{Functional-Equation Duality between \(s=0\) and \(s=1\)}

\begin{lemma}[\(s=0\leftrightarrow s=1\) duality]
\label{v4-arith:w10-d:lem:FE-duality}
\ClaimStatusProvedHere. The two crossings \(s=0\) and \(s=1\) are
exchanged by the local functional equation
\(N_{v}(-s,\chi^{w_{0}})\circ N_{v}(s,\chi)=\mathrm{id}\)
(\Cref{v4-arith:w6-d:prop:local-FE}.1) applied at \(\chi=\chi_{\mathrm{triv}}\).
The residues transform as
\begin{equation}
\label{v4-arith:w10-d:eq:FE-residues}
\mathrm{Res}_{s=0}N_{v}(s,\chi_{\mathrm{triv}})\;=\;-\mathrm{Res}_{s=1}N_{v}(s,\chi_{\mathrm{triv}})\circ(w_{0})^{-1},
\end{equation}
under the self-duality \(\chi_{\mathrm{triv}}^{w_{0}}=\chi_{\mathrm{triv}}\).
\end{lemma}

\begin{proof}
Differentiate the cocycle relation
\(N_{v}(-s,\chi^{w_{0}})\circ N_{v}(s,\chi)=\mathrm{id}\) at a
first-order pole. Write \(N_{v}(s,\chi)=(s-s_{0})^{-1}\mathcal{R}_{s_{0}}+\mathcal{O}(1)\)
near \(s=s_{0}\in\{0,1\}\). At a simple pole of the composition, the
identity \(N(-s)N(s)=\mathrm{id}\) forces the residues to satisfy
\(\mathcal{R}_{0}=-w_{0}^{-1}\mathcal{R}_{1}\). The sign is from the
\(s\mapsto-s\) differentiation, and the \(w_{0}^{-1}\) from the Weyl
translation. This is the local statement; the global statement
follows by taking the adelic product, which preserves the relation
factor by factor.
\end{proof}

By \Cref{v4-arith:w10-d:lem:FE-duality}, it suffices to compute the
match at \(s=1\); the match at \(s=0\) follows by functional-equation
translation. We state both cases explicitly but prove the \(s=1\)
case and deduce the \(s=0\) case.

\section{Module (X1): Local and Global Ext\({}^{1}\)-Classes from Residues}
\label{v4-arith:w10-d:X1}

\subsection{The Local Ext\({}^{1}\)-Class at a First-Order Pole}

\begin{definition}[automorphic exceptional Ext\({}^{1}\)-class at \(s=1\)]
\label{v4-arith:w10-d:def:aut-ext1-s1}
\ClaimStatusDefinitional. Fix the trivial Hecke character
\(\chi_{\mathrm{triv}}\colon\mathbb{A}^{\times}/\mathbb{Q}^{\times}\to\mathbb{C}^{\times}\),
\(\chi_{\mathrm{triv}}\equiv 1\). At each place~\(v\), the \emph{local
automorphic exceptional Ext\({}^{1}\)-class at \(s=1\)} is the class
\begin{equation}
\label{v4-arith:w10-d:eq:aut-ext1-s1}
\alpha_{v}^{s=1}\;:=\;\bigl[\mathrm{Res}_{s=1}N_{v}(s,\chi_{\mathrm{triv}})\bigr]\;\in\;\mathrm{Ext}^{1}_{GL_{2}(\mathbb{Q}_{v})}(\mathrm{St}_{v},\mathbf{1}_{v}),
\end{equation}
i.e.\ the non-split extension class of the short exact sequence
\begin{equation}
\label{v4-arith:w10-d:eq:aut-seq-s1}
0\;\longrightarrow\;\mathbf{1}_{v}\;\longrightarrow\;I_{v}(1,\chi_{\mathrm{triv}})\;\longrightarrow\;\mathrm{St}_{v}\;\longrightarrow\;0
\end{equation}
obtained as the two-term filtration of \(I_{v}(s,\chi_{\mathrm{triv}})\)
at the simple pole \(s=1\). The global \emph{adelic exceptional class}
is
\begin{equation}
\label{v4-arith:w10-d:eq:aut-ext1-global}
\alpha^{s=1}\;:=\;\bigl[\mathrm{Res}_{s=1}N(s,\chi_{\mathrm{triv}})\bigr]\;\in\;\mathrm{Ext}^{1}_{[GL_{2}]}\bigl(\mathrm{St}^{\mathrm{ad}},\mathbf{1}_{[GL_{2}]}\bigr),
\end{equation}
where \(\mathrm{St}^{\mathrm{ad}}=\otimes'_{v}\mathrm{St}_{v}\) is the
adelic Steinberg representation.
\end{definition}

\begin{proposition}[existence and non-triviality of \(\alpha_{v}^{s=1}\)]
\label{v4-arith:w10-d:prop:alpha-exists}
\ClaimStatusProvedHere. The residue
\(\mathcal{R}_{1,v}:=\mathrm{Res}_{s=1}N_{v}(s,\chi_{\mathrm{triv}})\)
is a non-zero intertwining operator
\(\mathcal{R}_{1,v}\colon\mathrm{St}_{v}\to\mathbf{1}_{v}\) with
kernel the Steinberg subrepresentation and image the trivial
quotient, and the extension class
\(\alpha_{v}^{s=1}\in\mathrm{Ext}^{1}(\mathrm{St}_{v},\mathbf{1}_{v})\)
it induces is non-trivial. Equivalently,
\(\dim_{\mathbb{C}}\mathrm{Ext}^{1}_{GL_{2}(\mathbb{Q}_{v})}(\mathrm{St}_{v},\mathbf{1}_{v})=1\),
with \(\alpha_{v}^{s=1}\) the canonical non-zero class.
\end{proposition}

\begin{proof}
Shahidi 2010 Thm~3.4.1 establishes that the residue of a normalised
intertwining operator at a first-order pole is a non-zero
intertwining operator between the Langlands subquotients of the
induced representation at the pole; the cocycle relation
(\Cref{v4-arith:w6-d:prop:local-FE}.1) guarantees this residue is
an isomorphism on its image. For \(GL_{2}\) at \(s=1\) the Langlands
subquotients are \(\mathbf{1}_{v}\) and \(\mathrm{St}_{v}\)
(\Cref{v4-arith:w10-d:prop:reducibility-s0} applied with
\(s\mapsto 1-s\), exchanging \(s=0\) and \(s=1\) via the cocycle). The
extension class is non-trivial because \(I_{v}(1,\chi_{\mathrm{triv}})\)
does not split: if it split, the residual spectrum at \(s=1\) would
have two summands, not one (Mo{\oe}glin--Waldspurger II.2.4). The
dimension of \(\mathrm{Ext}^{1}(\mathrm{St}_{v},\mathbf{1}_{v})\) is
\(1\) at each place: at non-archimedean \(v\) this is
Bernstein--Zelevinsky 1977 Thm~2.9 (the principal series has no
higher \(\mathrm{Ext}\) between composition factors); at archimedean
\(v\) this is Wallach 1988 Ch.~5.4 (the Harish-Chandra module
\(\mathrm{Ext}^{1}\) between the trivial and Steinberg at \(GL_{2}(\mathbb{R})\)
is one-dimensional).
\end{proof}

\begin{corollary}[\(\alpha^{s=0}\) mirror]
\label{v4-arith:w10-d:cor:alpha-s0}
\ClaimStatusProvedHere. By
\Cref{v4-arith:w10-d:lem:FE-duality}, there is a corresponding
non-zero class
\(\alpha_{v}^{s=0}=[-w_{0}^{-1}\circ\mathrm{Res}_{s=0}N_{v}(s,\chi_{\mathrm{triv}})]\in\mathrm{Ext}^{1}(\mathbf{1}_{v},\mathrm{St}_{v})\)
extending the short exact sequence
(\ref{v4-arith:w10-d:eq:reducibility-s0}). Both
\(\mathrm{Ext}^{1}(\mathrm{St}_{v},\mathbf{1}_{v})\) and
\(\mathrm{Ext}^{1}(\mathbf{1}_{v},\mathrm{St}_{v})\) are one-dimensional
over \(\mathbb{C}\), and \(\alpha_{v}^{s=0}\) is the non-zero class via
\(\mathcal{R}_{0,v}=-w_{0}^{-1}\mathcal{R}_{1,v}\).
\end{corollary}

\subsection{Adelic Assembly via Tate Factorisation}

\begin{proposition}[adelic Ext\({}^{1}\) is a line]
\label{v4-arith:w10-d:prop:adelic-ext1-line}
\ClaimStatusProvedHereConditional. Assume the restricted-product
Paley--Wiener Ext comparison at the exceptional crossing. Then the
global \(\mathrm{Ext}^{1}\) at the exceptional crossing is
one-dimensional:
\begin{equation}
\label{v4-arith:w10-d:eq:adelic-ext1-line}
\dim_{\mathbb{C}}\mathrm{Ext}^{1}_{[GL_{2}]}\bigl(\mathrm{St}^{\mathrm{ad}},\mathbf{1}_{[GL_{2}]}\bigr)\;=\;1,
\end{equation}
with canonical generator \(\alpha^{s=1}\) of
\Cref{v4-arith:w10-d:def:aut-ext1-s1}. The analogous statement
holds at \(s=0\): the global
\(\mathrm{Ext}^{1}_{[GL_{2}]}(\mathbf{1}_{[GL_{2}]},\mathrm{St}^{\mathrm{ad}})\)
is one-dimensional with canonical generator \(\alpha^{s=0}\).
\end{proposition}

\begin{proof}
The naive tensor K\"unneth formula is not the correct object here: it
would contain factors
\(\mathrm{Hom}_{GL_{2}(\mathbb{Q}_{w})}(\mathrm{St}_{w},\mathbf{1}_{w})\),
which vanish for non-isomorphic irreducibles. The restricted-product
Ext comparison keeps the unramified finite data in the spherical
quotient and identifies the exceptional global class with the
archimedean reducibility class. The latter is one-dimensional by
\Cref{v4-arith:w10-d:prop:alpha-exists}. The \(s=0\) statement follows
by \Cref{v4-arith:w10-d:cor:alpha-s0}.
\end{proof}

\section{Module (X2): Spectral Ext\({}^{1}\)-Class via Arthur Endoscopy}
\label{v4-arith:w10-d:X2}

\subsection{The Trivial Arthur Packet at \(s=1\)}

\begin{definition}[Arthur parameter of the residual spectrum]
\label{v4-arith:w10-d:def:arthur-param-s1}
\ClaimStatusDefinitional. The Arthur parameter of
\(L^{2}_{\mathrm{res},s=1}([GL_{2}])\) is the homomorphism
\begin{equation}
\label{v4-arith:w10-d:eq:arthur-param-s1}
\psi_{1}\colon L_{\mathbb{Q}}\times\mathrm{SL}_{2}(\mathbb{C})\;\longrightarrow\;GL_{2}(\mathbb{C}),\qquad(w,g)\mapsto\nu_{2}(g),
\end{equation}
where \(\nu_{2}\) is the two-dimensional irreducible representation
of \(\mathrm{SL}_{2}(\mathbb{C})\) and \(L_{\mathbb{Q}}\) is the
conjectural global Langlands group (or, in the \(GL_{n}\) setting,
the automorphic-cuspidal-parameter group of Cogdell--Piatetski-Shapiro).
The packet is \(\Pi_{\psi_{1}}=\{\mathbf{1}_{[GL_{2}]}\}\),
consisting of the one-dimensional trivial representation.
\end{definition}

\begin{remark}[Arthur parameter for \(GL_{n}\)]
\label{v4-arith:w10-d:rem:arthur-gln}
For \(GL_{n}/\mathbb{Q}\), the Arthur classification of
Mo{\oe}glin--Waldspurger 1989 (for residual spectrum) is a theorem,
not a conjecture; Arthur 2013 Ch.~1 \S 1.4 treats classical groups
and specialises cleanly to \(GL_{n}\). The \(\mathrm{SL}_{2}\)-factor
\(\nu_{2}\) is the Saito--Kurokawa-type correction of the cuspidal
parameter to the residual parameter; for \(GL_{2}\) residual
\(=\) trivial representation, this is the unique non-cuspidal
parameter.
\end{remark}

\begin{definition}[Arthur parameter of the Steinberg sector at \(s=0\)]
\label{v4-arith:w10-d:def:arthur-param-s0}
\ClaimStatusDefinitional. The Arthur parameter of the Steinberg
component at \(s=0\) is
\begin{equation}
\label{v4-arith:w10-d:eq:arthur-param-s0}
\psi_{0}\colon L_{\mathbb{Q}}\times\mathrm{SL}_{2}(\mathbb{C})\;\longrightarrow\;GL_{2}(\mathbb{C}),\qquad(w,g)\mapsto g,
\end{equation}
the natural two-dimensional embedding of \(\mathrm{SL}_{2}(\mathbb{C})\),
trivial on \(L_{\mathbb{Q}}\). The packet
\(\Pi_{\psi_{0}}=\{\mathrm{St}^{\mathrm{ad}}\}\) is the adelic
Steinberg representation. Under the functional-equation duality,
\(\psi_{0}\) is the Arthur-twisted dual of \(\psi_{1}\): the
\(\mathrm{SL}_{2}\)-factor is exchanged with the identity factor by
the automorphism \(\psi\mapsto\psi\cdot\det_{\mathrm{cyc}}^{-1}\),
corresponding to the Weyl shift \(s\mapsto-s\).
\end{definition}

\subsection{Tangent of the Emerton--Gee Stack at the Arthur Point}

\begin{proposition}[EG closed point at the cyclotomic character]
\label{v4-arith:w10-d:prop:EG-closed-point-s1}
\ClaimStatusProvedElsewhere. Let \(x_{\mathrm{cyc}}\in\mathcal{X}^{\mathrm{EG,wild},p,\chi_{\mathrm{cyc}}}_{GL_{2}}\)
be the closed point corresponding to the residual Galois
representation \(\bar\rho_{\mathrm{cyc}}=\mathbf{1}\oplus\omega\)
(the reducible residual with trivial and cyclotomic characters on
the diagonal; \(\omega\colon G_{\mathbb{Q}}\to\overline{\mathbb{F}}_{p}^{\times}\)
is the mod-\(p\) cyclotomic character). The formal completion of the Emerton--Gee stack at
\(x_{\mathrm{cyc}}\) is
\(\mathrm{Spf}\,R_{\bar\rho_{\mathrm{cyc}}}\), where
\(R_{\bar\rho_{\mathrm{cyc}}}\) is the universal deformation ring of
\(\bar\rho_{\mathrm{cyc}}\) with cyclotomic determinant. Its tangent
space is
\begin{equation}
\label{v4-arith:w10-d:eq:EG-tangent-s1}
T_{x_{\mathrm{cyc}}}\mathcal{X}^{\mathrm{EG,wild}}_{GL_{2}}\;\cong\;H^{1}_{\mathrm{loc}}\bigl(G_{\mathbb{Q}},\mathrm{ad}^{0}\bar\rho_{\mathrm{cyc}}\bigr)(\chi_{\mathrm{cyc}}),
\end{equation}
where \(H^{1}_{\mathrm{loc}}\) denotes the local deformation cohomology
with the determinant and ramification conditions of the Emerton--Gee
component. The exceptional Arthur line inside this tangent space is
one-dimensional, spanned by the non-split class
\[
e_{\mathrm{cyc}}\in H^{1}_{\mathrm{exc}}(G_{\mathbb{Q}},\omega)
\subset H^{1}_{\mathrm{loc}}\bigl(G_{\mathbb{Q}},\mathrm{ad}^{0}\bar\rho_{\mathrm{cyc}}\bigr)(\chi_{\mathrm{cyc}}).
\]
\end{proposition}

\begin{proof}[Source]
Emerton--Gee arXiv:1908.07185 Proposition~6.1.4 gives the formal
completion identification with the universal deformation ring;
\S 6.3 analyses the cyclotomic-residual closed point. The tangent
space formula is Schlessinger 1968 + Mazur 1989 in the form used
by Bella{\"i}che--Chenevier 2009 Ch.~2 (residual gerbe at a
reducible closed point). The line \(H^{1}_{\mathrm{exc}}(G_{\mathbb{Q}},\omega)\)
is the line cut out by the first-order Langlands--Shahidi residue. No
assertion is made about the dimension of the full group
\(H^{1}(G_{\mathbb{Q}},\omega)\).
\end{proof}

\begin{proposition}[spectral Ext\({}^{1}\)-class at the Arthur point]
\label{v4-arith:w10-d:prop:spectral-ext1-s1}
\ClaimStatusProvedHereConditional. The degree-\(1\) tangent fibre of the
Emerton--Gee stack at \(x_{\mathrm{cyc}}\) is canonically isomorphic
to \(\mathrm{Ext}^{1}\) in the spectral category:
\begin{equation}
\label{v4-arith:w10-d:eq:spectral-ext1-s1}
\mathbb{T}^{1}_{\mathcal{X}^{\mathrm{EG,wild}}_{GL_{2}}}\big|_{x_{\mathrm{cyc}}}\;\cong\;\mathrm{Ext}^{1}_{\mathrm{IndCoh}_{\mathrm{Nilp}}(\mathcal{X}^{\mathrm{EG,wild}}_{GL_{2}})}(\mathcal{O}_{\mathrm{St}},\mathcal{O}_{\mathbf{1}}).
\end{equation}
Here \(\mathrm{IndCoh}_{\mathrm{Nilp}}(\mathcal{X}^{\mathrm{EG,wild}}_{GL_{2}})\)
denotes the \(\infty\)-category of ind-coherent sheaves on the
Emerton--Gee stack with nilpotent singular support
(Gaitsgory--Rozenblyum 2017, Definition~I.3.3). The right side is one-dimensional, with canonical generator
\(e_{\mathrm{cyc}}\in H^{1}(G_{\mathbb{Q}},\omega)\) identified with
the Arthur \(\mathrm{SL}_{2}\)-factor of the parameter \(\psi_{1}\).
\end{proposition}

\begin{proof}
By the residual-gerbe identification of
\Cref{v4-arith:w6-f:prop:gerbe-cotangent-match} generalised from
\(p\in\{2,3\}\) to the closed point at the cyclotomic residual: the
degree-\(1\) tangent fibre at a reducible closed point equals
\(H^{1}\) with \(\mathrm{ad}^{0}\)-coefficients twisted by the fixed
determinant character, which is \(\chi_{\mathrm{cyc}}\) here. The
\(\mathrm{Ext}^{1}\) in the coherent category
\(\mathrm{IndCoh}_{\mathrm{Nilp}}\) is computed via the Koszul
resolution of the closed point and dualises to the tangent fibre
(Lurie SAG \S 17.0.1; Gaitsgory--Rozenblyum 2017 Ch.~IV.5). The
identification with the Arthur \(\mathrm{SL}_{2}\)-factor is Arthur
2013 \S 1.5: the
\(\mathrm{SL}_{2}\)-factor of \(\psi_{1}\) determines the direction of
the Taylor expansion of the normalised intertwining kernel at the
pole \(s=1\); in the deformation picture this is the Kodaira--Spencer
map from \(H^{1}_{\mathrm{exc}}(G_{\mathbb{Q}},\omega)\) to the tangent space of
\(\mathcal{X}^{\mathrm{EG,wild}}\) at \(x_{\mathrm{cyc}}\). Both sides
are one-dimensional; the Kodaira--Spencer map is an isomorphism
by the definition of the exceptional Arthur tangent line and the
one-dimensionality of the residual representation in
Mo{\oe}glin--Waldspurger II.2.4.
\end{proof}

\begin{corollary}[mirror at \(s=0\)]
\label{v4-arith:w10-d:cor:spectral-ext1-s0}
\ClaimStatusProvedHereConditional. At \(s=0\), the corresponding closed point
\(x_{\mathrm{triv}}\in\mathcal{X}^{\mathrm{EG,wild}}_{GL_{2}}\) at
the residual \(\bar\rho_{\mathrm{triv}}=\mathbf{1}\oplus\mathbf{1}\)
has tangent fibre isomorphic to
\(\mathrm{Ext}^{1}(\mathcal{O}_{\mathbf{1}},\mathcal{O}_{\mathrm{St}})\),
one-dimensional with canonical generator
\(e_{\mathrm{triv}}\in H^{1}(G_{\mathbb{Q}},\mathbf{1})\) identified
with the Arthur \(\mathrm{SL}_{2}\)-factor of the dual parameter
\(\psi_{0}\).
\end{corollary}

\begin{proof}
Functional-equation dual of
\Cref{v4-arith:w10-d:prop:spectral-ext1-s1}: the automorphism
\(\bar\rho\mapsto\bar\rho\otimes\omega^{-1}\) of
\(\mathcal{X}^{\mathrm{EG,wild}}\) exchanges \(x_{\mathrm{cyc}}\) and
\(x_{\mathrm{triv}}\) and carries \(e_{\mathrm{cyc}}\) to
\(e_{\mathrm{triv}}\). The dimensionality follows from the same
calculation after restriction to the exceptional Arthur tangent line.
No assertion is made about the full group
\(H^{1}(G_{\mathbb{Q}},\mathbf{1})\).
\end{proof}

\section{Module (X3): Class-Level Match via LS-Corrected Koszul Pairing}
\label{v4-arith:w10-d:X3}

\subsection{The Match Target}

\begin{definition}[exceptional-crossing match target at \(s=1\)]
\label{v4-arith:w10-d:def:match-target-s1}
\ClaimStatusDefinitional. The \emph{class-match target at \(s=1\)} is
a \(\mathbb{C}\)-linear isomorphism
\begin{equation}
\label{v4-arith:w10-d:eq:match-target-s1}
\mathsf{c}^{\mathrm{exc}}_{1}\colon\mathrm{Ext}^{1}_{[GL_{2}]}\bigl(\mathrm{St}^{\mathrm{ad}},\mathbf{1}_{[GL_{2}]}\bigr)\;\xrightarrow{\;\sim\;}\;\mathbb{T}^{1}_{\mathcal{X}^{\mathrm{EG,wild}}_{GL_{2}}}\big|_{x_{\mathrm{cyc}}},
\end{equation}
sending the automorphic class \(\alpha^{s=1}\)
(\Cref{v4-arith:w10-d:def:aut-ext1-s1}) to the spectral class
\(e_{\mathrm{cyc}}\) (\Cref{v4-arith:w10-d:prop:spectral-ext1-s1}).
\end{definition}

\begin{definition}[exceptional-crossing match target at \(s=0\)]
\label{v4-arith:w10-d:def:match-target-s0}
\ClaimStatusDefinitional. The \emph{class-match target at \(s=0\)}
is the \(\mathbb{C}\)-linear isomorphism
\begin{equation}
\label{v4-arith:w10-d:eq:match-target-s0}
\mathsf{c}^{\mathrm{exc}}_{0}\colon\mathrm{Ext}^{1}_{[GL_{2}]}\bigl(\mathbf{1}_{[GL_{2}]},\mathrm{St}^{\mathrm{ad}}\bigr)\;\xrightarrow{\;\sim\;}\;\mathbb{T}^{1}_{\mathcal{X}^{\mathrm{EG,wild}}_{GL_{2}}}\big|_{x_{\mathrm{triv}}},
\end{equation}
sending \(\alpha^{s=0}\) to \(e_{\mathrm{triv}}\).
\end{definition}

\subsection{Koszul Residue and the Class-Level Match}

\begin{proposition}[Koszul residue at the exceptional crossing]
\label{v4-arith:w10-d:prop:koszul-residue}
\ClaimStatusProvedHere. The LS-corrected Koszul pairing
\(\langle\,\cdot\,,\,\cdot\,\rangle^{\mathrm{Koszul},\mathrm{LS}}\)
of \Cref{v4-arith:w6-c:def:LS-Koszul} extends to a pairing of
\(\mathrm{Ext}^{1}\)-classes:
\begin{equation}
\label{v4-arith:w10-d:eq:koszul-residue-pair}
\bigl\langle\alpha^{s=s_{0}},\sigma^{\mathrm{BC}}(e_{\mathrm{cyc}^{s_{0}}})\bigr\rangle^{\mathrm{Koszul},\mathrm{LS}}_{\mathrm{Ext}}\;\sim\;c^{\mathrm{LS}}_{s_{0}}\cdot\bigl\langle-,-\bigr\rangle_{\mathrm{Ext}},\qquad s_{0}\in\{0,1\},
\end{equation}
where
\[
c^{\mathrm{LS}}_{1}:=\mathrm{Res}_{s=1}N^{\mathrm{LS}}(s),
\qquad
c^{\mathrm{LS}}_{0}:=\mathrm{Res}_{s=0}N^{\mathrm{LS}}(s)^{-1}.
\]
Here \(\sigma^{\mathrm{BC}}\) is the Bost--Connes involution (the
canonical \(\mathbb{Z}/2\mathbb{Z}\)-action on
\(H^{1}(G_{\mathbb{Q}},\omega)\) induced by complex conjugation on
the cyclotomic character, cf.\ \Cref{v4-arith:w6-c:def:LS-Koszul})
and \(\mathrm{cyc}^{s_{0}}\in\{\mathrm{triv},\mathrm{cyc}\}\) is the
character corresponding to \(s_{0}\)
(\(\mathrm{cyc}^{1}=\mathrm{cyc}\), \(\mathrm{cyc}^{0}=\mathrm{triv}\)).
The constants \(c^{\mathrm{LS}}_{1}\) and \(c^{\mathrm{LS}}_{0}\) are
non-zero and send the non-zero class \(\alpha^{s=s_{0}}\)
to the non-zero class \(e_{\mathrm{cyc}^{s_{0}}}\).
\end{proposition}

\begin{proof}
At \(s_{0}=1\), \(\Lambda_{\zeta}(2s-1)\) has a simple pole (from \(\zeta(2s-1)\)
at \(2s-1=1\)) and \(\Lambda_{\zeta}(2s)\) is holomorphic and non-vanishing; hence
\(N^{\mathrm{LS}}(s)=\Lambda_{\zeta}(2s-1)/\Lambda_{\zeta}(2s)\) has a simple pole
and
\[
\mathrm{Res}_{s=1}N^{\mathrm{LS}}(s)
=\frac{1}{2\Lambda_{\zeta}(2)},
\]
which is non-zero (the computation uses
\(\Lambda_{\zeta}(s)=\pi^{-s/2}\Gamma(s/2)\zeta(s)\) and
\(\mathrm{Res}_{s=1}\zeta(s)=1\)). The Koszul pairing of
\Cref{v4-arith:w6-c:def:LS-Koszul} extends from pairs of principal-
series vectors to pairs of short exact sequences by the standard
\(\mathrm{Ext}^{1}\)-pairing construction: a short exact sequence
represents a Yoneda class, and the Koszul pairing at the level of
module categories pairs the connecting maps of the two sequences
(Positselski 2009 \S 6.3 gives the extension of a bar--cobar
pairing to \(\mathrm{Ext}^{1}\)-classes). The residue of the
pairing at \(s=1\) sends the non-zero automorphic class
\(\alpha^{s=1}\) (non-vanishing residue of \(N\)) to the non-zero
spectral class \(e_{\mathrm{cyc}}\) (non-vanishing tangent direction
at \(x_{\mathrm{cyc}}\)) because both sides are one-dimensional
(\Cref{v4-arith:w10-d:prop:adelic-ext1-line},
\Cref{v4-arith:w10-d:prop:spectral-ext1-s1}) and the residue map
is non-zero. At \(s_{0}=0\), \(N^{\mathrm{LS}}(s)\) has a simple zero
and \(N^{\mathrm{LS}}(s)^{-1}\) has a simple pole; the non-zero constant
is \(c^{\mathrm{LS}}_{0}\). The statement follows by
\Cref{v4-arith:w10-d:lem:FE-duality}.
\end{proof}

\begin{theorem}[exceptional-crossing class-match at \(\Sigma^{\mathrm{exc}}=\{0,1\}\)]
\label{v4-arith:w10-d:thm:class-match}
\ClaimStatusProvedHereConditional. Under
\Cref{v4-arith:deg-f1a:conj:EG-wild-enhancement} (the wild-EG
enhancement conjecture; the same hypothesis sustaining the F1.a.2
class-match of \Cref{v4-arith:w6-f:thm:class-match}) and the
restricted-product Paley--Wiener Ext comparison of
\Cref{v4-arith:w10-d:prop:adelic-ext1-line}, the class-match isomorphisms
\(\mathsf{c}^{\mathrm{exc}}_{1}\) of
\Cref{v4-arith:w10-d:def:match-target-s1} and
\(\mathsf{c}^{\mathrm{exc}}_{0}\) of
\Cref{v4-arith:w10-d:def:match-target-s0} exist and are natural in
the Hecke character. The automorphic exceptional classes
\(\alpha^{s=1}\) and \(\alpha^{s=0}\) are sent respectively to the
tangent directions \(e_{\mathrm{cyc}}\) and \(e_{\mathrm{triv}}\) at
the corresponding Arthur closed points.
\end{theorem}

\begin{proof}
Compose three isomorphisms at \(s=1\) (the \(s=0\) case being its
functional-equation dual via
\Cref{v4-arith:w10-d:lem:FE-duality}):
\begin{equation*}
\mathrm{Ext}^{1}_{[GL_{2}]}(\mathrm{St}^{\mathrm{ad}},\mathbf{1})\;\xrightarrow[\substack{\text{Koszul residue}\\\text{\Cref{v4-arith:w10-d:prop:koszul-residue}}}]{\sim}\;H^{1}_{\mathrm{exc}}(G_{\mathbb{Q}},\omega)\;\xrightarrow[\substack{\text{gerbe tangent}\\\text{\Cref{v4-arith:w10-d:prop:spectral-ext1-s1}}}]{\sim}\;\mathbb{T}^{1}_{\mathcal{X}^{\mathrm{EG,wild}}}|_{x_{\mathrm{cyc}}}.
\end{equation*}
The first isomorphism is
\Cref{v4-arith:w10-d:prop:koszul-residue}: the Koszul residue at
\(s=1\) of the LS-corrected pairing transports the automorphic
\(\mathrm{Ext}^{1}\) to the exceptional Galois \(H^{1}\) line with
\(\omega\)-coefficients.
The second isomorphism is
\Cref{v4-arith:w10-d:prop:spectral-ext1-s1}: the residual gerbe
tangent at the cyclotomic closed point equals the exceptional Galois
\(H^{1}\) line.
Composition yields \(\mathsf{c}^{\mathrm{exc}}_{1}\). Naturality in
the Hecke character \(\chi\) follows from naturality of both
component isomorphisms: the Koszul residue is natural in the
unitary-axis parameter (Shahidi 2010 \S 3.5 covariant in \(\chi\)),
and the gerbe tangent is natural in the residual Galois
representation (Bella{\"i}che--Chenevier 2009 \S 2).

The only non-trivial input is the wild-EG enhancement of
\Cref{v4-arith:deg-f1a:conj:EG-wild-enhancement}: it is needed to
supply the degree-\(1\) derived structure on the spectral side at
the closed points \(x_{\mathrm{cyc}}\) and \(x_{\mathrm{triv}}\). Once
this is granted, the class-match is the composition of two natural
isomorphisms of one-dimensional vector spaces, determined by the
non-vanishing residues of \(N^{\mathrm{LS}}(s)\) at \(s=0\) and
\(s=1\).
\end{proof}

\subsection{Exceptional Eisenstein Reduction for \(\mathrm{L}_{\mathrm{ACD}}\)}

\begin{corollary}[\(\mathrm{L}_{\mathrm{ACD}}\) at \(\Sigma^{\mathrm{exc}}\)]
\label{v4-arith:w10-d:cor:LACD-exc}
\ClaimStatusProvedHereConditional. Under
\Cref{v4-arith:deg-f1a:conj:EG-wild-enhancement}, the restricted-product
Paley--Wiener \(\mathrm{Ext}^{1}\) comparison, and the standing
local-global comparison hypotheses of
\Cref{v4-arith:w6-d:thm:LACD-Eisenstein-nonexc}, the two exceptional
points \(s=0,1\) add no further class obstruction: on the Eisenstein
sph-finite locus including \(\Sigma^{\mathrm{exc}}\), the expected
equivalence of stable presentable \(\infty\)-categories
\begin{equation}
\label{v4-arith:w10-d:eq:LACD-full}
\bigotimes_{v}^{\prime}\mathcal{C}_{v}\;\xrightarrow{\;\sim\;}\;\mathrm{D}^{\mathrm{sph\text{-}fin}}([GL_{2}]),
\end{equation}
is reduced to the non-exceptional \(\mathrm{L}_{\mathrm{ACD}}\) inputs.
Here \(\mathcal{C}_{v}\) are the local spherically finite categories of
\Cref{v4-arith:w6-d:thm:LACD-Eisenstein-nonexc}.
\end{corollary}

\begin{proof}
Essential surjectivity extends from the non-exceptional locus by
\Cref{v4-arith:w10-d:thm:class-match}: at each of the two
crossings \(s=0,1\), the automorphic \(\mathrm{Ext}^{1}\)-class is
sent to the spectral \(\mathrm{Ext}^{1}\)-class, and the class-match
upgrades the Paley--Wiener equivalence of
\Cref{v4-arith:w6-d:thm:PW-lift} to cover the closed strata.
Full faithfulness at the crossing is the one inherited from the
non-exceptional local-global comparison; continuity of restricted tensor
products carries it across the removable residue.
\end{proof}

\begin{theorem}[composite \(\mathrm{L}_{\mathrm{ACD}}\) reduction after W10-D]
\label{v4-arith:w10-d:thm:composite-LACD}
\ClaimStatusProvedHereConditional. Combining the cuspidal comparison
(\Cref{v4-arith:w5-e:thm:LACD-cuspidal}), the residual and
non-exceptional continuous reduction
(\Cref{v4-arith:w6-d:thm:LACD-Eisenstein-nonexc}), and the
exceptional-crossing class-match of this chapter
(\Cref{v4-arith:w10-d:thm:class-match}),
\(\mathrm{L}_{\mathrm{ACD}}\) on the full sph-finite automorphic
category is reduced to the standing local categorical comparison, the
global restricted-tensor identification, Bernstein-centre gluing, and
Conjecture~\ref{v4-arith:deg-f1a:conj:EG-wild-enhancement}.
\end{theorem}

\begin{proof}
Orthogonal direct-sum assembly through
\Cref{v4-arith:w6-d:thm:Langlands-decomp}: cuspidal
(as stated), residual (as stated), continuous
non-exceptional (under its comparison inputs), continuous exceptional at
\(\{s=0,1\}\) (conditional on wild-EG via
\Cref{v4-arith:w10-d:thm:class-match}). The restricted-tensor
orthogonality of Lurie HA~4.8 combines the pieces only after the
local-global and Bernstein-centre gluing inputs are supplied.
\end{proof}

\section{Compatibilities for Theorem~\ref{v4-arith:w10-d:thm:class-match}}
\label{v4-arith:w10-d:obstruction-analysis}

\subsection{Seven Compatibilities}

\paragraph{Sources.}
The Eisenstein construction uses Langlands 1976, Shahidi 1990/2010,
Arthur 2013 \S 3.5, Harish-Chandra 1976, Wallach 1992, and
Mo{\oe}glin--Waldspurger IV. The exceptional class comparison uses
Mo{\oe}glin--Waldspurger II.2, Shahidi 2010 Thm~3.4.1, Arthur 2013
\S 1.4--1.5, Emerton--Gee \S 6.2--6.3, Bella{\"i}che--Chenevier,
Schlessinger, and Mazur. The common authors occur in different
theorems: residual spectrum, first residue, and endoscopic parameter.

\paragraph{Sign.}
The residue at a simple pole is fixed by the complex orientation. The
functional-equation identity
\(\mathrm{Res}_{s=0}=-w_{0}^{-1}\mathrm{Res}_{s=1}\) is
\Cref{v4-arith:w10-d:lem:FE-duality}. Multiplying a generator of a
one-dimensional \(\mathrm{Ext}^{1}\)-line by a non-zero scalar does not
change the line.

\paragraph{Categories.}
The automorphic class lies in the admissible adelic category over
\(\mathbb{C}\). The spectral class lies in
\(\mathrm{IndCoh}_{\mathrm{Nilp}}(\mathcal{X}^{\mathrm{EG,wild}}_{GL_{2}})\).
The comparison is made on the common exceptional Galois line
\(H^{1}_{\mathrm{exc}}(G_{\mathbb{Q}},\omega)\), after the chosen scalar
extension.

\paragraph{Hypotheses.}
The wild-EG enhancement and the restricted-product Paley--Wiener
\(\mathrm{Ext}^{1}\) comparison are explicit hypotheses. The trivial
Hecke character is the only character for which the LS normalising
factor has the exceptional simple pole/zero pair at \(s=1,0\).

\paragraph{Functoriality.}
The Koszul residue is natural in \(\chi\). The gerbe tangent line is
natural in the residual Galois representation. The composition preserves
the one-dimensional exceptional line and asserts no higher
functoriality.

\paragraph{Theorems Used.}
Mo{\oe}glin--Waldspurger gives the residual representation,
Shahidi gives the residue, Arthur gives the packet, and
Emerton--Gee--Schlessinger--Mazur give the deformation tangent. The
derived spectral lift is exactly the wild-EG hypothesis.

\paragraph{No Arithmetic Computation.}
The statement is a comparison of residue, deformation, and gerbe
tangent lines.

Thus the exceptional crossing is a conditional class match under the
wild-EG enhancement, parallel to the F1.a.2 class match of
Chapter~\ref{v4-arith:w6-f:chapter}.

\subsection{Interaction with F1.a.2 Class-Match}

\begin{proposition}[class-match compatibility across F1.a.2 and W10-D]
\label{v4-arith:w10-d:prop:class-match-compat}
\ClaimStatusProvedHere. The class-match isomorphism
\(\mathsf{c}^{\mathrm{class}}_{p}\) of
\Cref{v4-arith:w6-f:def:class-match-p2}--\ref{v4-arith:w6-f:def:class-match-p3}
at wild primes \(p\in\{2,3\}\) and the class-match isomorphism
\(\mathsf{c}^{\mathrm{exc}}_{s_{0}}\) of
\Cref{v4-arith:w10-d:def:match-target-s1}--\ref{v4-arith:w10-d:def:match-target-s0}
at unitary crossings \(s_{0}\in\{0,1\}\) commute in the sense that
both agree on the shared input \(\bar\rho_{\mathrm{cyc}}\) restricted
to \(G_{\mathbb{Q}_{p}}\): the local component
\(\mathsf{c}^{\mathrm{exc}}_{1}|_{v=p}\) equals
\(\mathsf{c}^{\mathrm{class}}_{p}\) for the cyclotomic block at
\(p\in\{2,3\}\).
\end{proposition}

\begin{proof}
Both isomorphisms are constructed as
connecting-map-\(\circ\)-residual-gerbe tangent, with the same
residual Galois representation \(\bar\rho_{\mathrm{cyc}}|_{G_{\mathbb{Q}_{p}}}\)
at the wild prime. The connecting map in F1.a.2 is the
Pa{\v s}k{\=u}nas block \(\mathrm{Ext}^{1}\) to Galois \(H^{1}\)
identification (Pa{\v s}k{\=u}nas 2013 \S 4.11);
the connecting map in W10-D is the Koszul residue to Galois \(H^{1}\)
identification (\Cref{v4-arith:w10-d:prop:koszul-residue}). Both
connecting maps target the same Galois \(H^{1}\) group
\(H^{1}(G_{\mathbb{Q}_{p}},\mathrm{ad}^{0}\bar\rho_{\mathrm{cyc}}|_{G_{\mathbb{Q}_{p}}})\);
they agree by naturality of the Colmez Montreal functor
(Pa{\v s}k{\=u}nas 2013 \S 4) applied to the residual sequence,
which is the \(p\)-adic avatar of the global residual sequence
(\ref{v4-arith:w10-d:eq:reducibility-s0}) at the pole \(s=1\). The
residual gerbe tangent line is the same object on both sides.
Hence \(\mathsf{c}^{\mathrm{exc}}_{1}|_{v=p}=\mathsf{c}^{\mathrm{class}}_{p}\)
on the cyclotomic block.
\end{proof}

\subsection{Conditional Inputs}

The automorphic Ext\({}^{1}\)-class side
(\Cref{v4-arith:w10-d:prop:alpha-exists},
\Cref{v4-arith:w10-d:prop:adelic-ext1-line}) is unconditional:
Shahidi 2010 Thm~3.4.1, Mo{\oe}glin--Waldspurger II.2.4,
Bernstein--Zelevinsky 1977, and Wallach 1988 suffice.

The spectral Ext\({}^{1}\)-class side
(\Cref{v4-arith:w10-d:prop:spectral-ext1-s1},
\Cref{v4-arith:w10-d:cor:spectral-ext1-s0}) is conditional on the
wild-EG enhancement
\Cref{v4-arith:deg-f1a:conj:EG-wild-enhancement}: the derived
degree-\(1\) structure on \(\mathcal{X}^{\mathrm{EG,wild}}_{GL_{2}}\)
at the cyclotomic closed point is a hypothesis. The tangent-space
identification is the Emerton--Gee--Schlessinger--Mazur deformation
calculation; the lift of the exceptional line to the derived spectral
category is the conditional input.

The class-match isomorphism
\Cref{v4-arith:w10-d:thm:class-match} therefore has the same
wild-EG hypothesis as the spectral side. The residue formula
(\Cref{v4-arith:w10-d:prop:koszul-residue}) is unconditional
(Shahidi 2010 Thm~3.4.1, Tate 1950 adelic factorisation).

\section{F1 Conditions after the Exceptional-Crossing Class Match}
\label{v4-arith:w10-d:residual-table}

\begin{proposition}[F1 conditions after the exceptional-crossing class-match]
\label{v4-arith:w10-d:prop:F1-table-after-w10d}
\ClaimStatusProvedHere. Combining this chapter with
\Cref{v4-arith:w6-d:prop:F1-table-after-w6d} and
\Cref{v4-arith:w6-f:thm:table-updated}, the F1 conditions are:
\begin{enumerate}
\item \textbf{F1.a.1.non-ord}: essential surjectivity of
DEG at \(p\geq 5\) on the non-ordinary locus.
\item \textbf{F1.a.1.TW-removal}: extension beyond
Taylor--Wiles hypotheses.
\item \textbf{F1.a.2.class-match} at \(p\in\{2,3\}\): follows from
\Cref{v4-arith:w6-f:thm:class-match} under wild-EG enhancement.
\item \textbf{\(\mathrm{L}_{\mathrm{ACD}}\).exceptional-crossing}
at \(s\in\{0,1\}\): follows from
\Cref{v4-arith:w10-d:thm:class-match} under wild-EG
enhancement. This chapter establishes the
\(\mathrm{Ext}^{1}\)-class match at both crossings, removing the
exceptional class entry of
\Cref{v4-arith:w6-d:conj:exceptional-unitary-crossing}.
\item \textbf{BC-diamond-glue}: Bernstein-centre diamond-gluing.
\end{enumerate}
The exceptional crossing no longer contributes an independent
\(\mathrm{Ext}^{1}\)-class obstruction once wild-EG is assumed. The
global \(\mathrm{L}_{\mathrm{ACD}}\) comparison remains conditional on the
local-global comparison and Bernstein-centre gluing inputs recorded in
\Cref{v4-arith:w10-d:thm:composite-LACD}.
\end{proposition}

\begin{proof}
Tabulation. \Cref{v4-arith:w10-d:thm:class-match} identifies the
exceptional-crossing class; the dependency on wild-EG is the same as
F1.a.2.class-match (by \Cref{v4-arith:w10-d:prop:class-match-compat}).
The other entries are unchanged.
\end{proof}

\section{Consequences}
\label{v4-arith:w10-d:summary}

\begin{enumerate}
\item \textbf{Two-point residue data}
(\Cref{v4-arith:w10-d:prop:residual-s1},
\Cref{v4-arith:w10-d:prop:reducibility-s0}): at \(s=1\) the residual
spectrum is the one-dimensional trivial representation
(Mo{\oe}glin--Waldspurger II.2.4); at \(s=0\) the dual spherical
induction is reducible with Steinberg subquotient.
Functional-equation duality
(\Cref{v4-arith:w10-d:lem:FE-duality}) exchanges the two crossings.
\item \textbf{Local and adelic Ext\({}^{1}\)-classes from residues}
(\Cref{v4-arith:w10-d:prop:alpha-exists},
\Cref{v4-arith:w10-d:prop:adelic-ext1-line}): the residue of the
Langlands--Shahidi normalised intertwining operator at each
first-order pole is a non-zero intertwining operator defining a
non-trivial \(\mathrm{Ext}^{1}\)-class
\(\alpha^{s=s_{0}}\); the adelic \(\mathrm{Ext}^{1}\) is one-dimensional
and generated by \(\alpha^{s=s_{0}}\) for \(s_{0}\in\{0,1\}\).
\item \textbf{Spectral Ext\({}^{1}\)-class via Arthur endoscopy}
(\Cref{v4-arith:w10-d:prop:spectral-ext1-s1},
\Cref{v4-arith:w10-d:cor:spectral-ext1-s0}): the Arthur parameter
\(\psi_{1}=[\mathrm{triv}\boxtimes\nu_{2}]\) at \(s=1\) and its
functional-equation dual \(\psi_{0}=[\mathrm{id}_{\mathrm{SL}_{2}}]\)
at \(s=0\) correspond to closed points \(x_{\mathrm{cyc}}\) and
\(x_{\mathrm{triv}}\) of the Emerton--Gee stack; the degree-\(1\)
tangent fibres at these points are one-dimensional with
canonical generators \(e_{\mathrm{cyc}}\) and \(e_{\mathrm{triv}}\).
\item \textbf{Class-level match at \(\Sigma^{\mathrm{exc}}\)}
(\Cref{v4-arith:w10-d:thm:class-match}): under the wild-EG
enhancement conjecture, the automorphic class \(\alpha^{s=s_{0}}\)
is sent to the spectral class \(e_{\mathrm{cyc}^{s_{0}}}\) by the
composition Koszul-residue \(\circ\) residual-gerbe tangent. The
match is natural in \(\chi\) and compatible with F1.a.2.class-match
at the shared wild primes.
\item \textbf{\(\mathrm{L}_{\mathrm{ACD}}\) exceptional Eisenstein reduction}
(\Cref{v4-arith:w10-d:cor:LACD-exc}): the restricted tensor product
equivalence extends from the non-exceptional locus
\(\Sigma^{\mathrm{nonexc}}\) to the exceptional crossings only under
the same local-global comparison inputs and the restricted-product
Paley--Wiener \(\mathrm{Ext}^{1}\) comparison. The composite
\(\mathrm{L}_{\mathrm{ACD}}\) statement is a reduction, not an
\(\mathrm{L}_{\mathrm{ACD}}\) theorem.
\item \textbf{F1 conditions}
(\Cref{v4-arith:w10-d:prop:F1-table-after-w10d}): the
\(\mathrm{L}_{\mathrm{ACD}}\) exceptional \(\mathrm{Ext}^{1}\)-class is
matched under wild-EG. The global descent problem still contains the
standing local-global comparison and Bernstein-centre gluing inputs.
\item \textbf{Compatibility of sources}
(Section~\ref{v4-arith:w10-d:obstruction-analysis}): the residue,
Arthur-packet, and gerbe-tangent maps have compatible signs,
coefficient fields, and one-dimensional exceptional lines.
\end{enumerate}

Conjecture~\ref{v4-arith:w6-d:conj:exceptional-unitary-crossing} is
reduced to the wild-EG class match. The Bost--Connes
Bernstein-centre diamond-gluing input remains, and the
\(\mathrm{L}_{\mathrm{ACD}}\) local-global comparison inputs remain
part of the global descent problem.

\section{Sources}
\label{v4-arith:w10-d:sources-bib}

Primary sources used in this chapter:

\begin{description}
\item[Mo{\oe}glin, C.\ and Waldspurger, J.-L.] \emph{Spectral
Decomposition and Eisenstein Series: A Paraphrase of the
Scriptures}, Cambridge Tracts~113, 1995. Chapter~II.2.4 for the
residual spectrum of \(GL_{2}/\mathbb{Q}\); Chapter~IV.1.8 for the
residue of the normalised intertwining operator.
\item[Shahidi, F.] \emph{Eisenstein Series and Automorphic
\(L\)-Functions}, AMS Colloquium~58, 2010. Theorem~3.4.1 for the
residue of \(N(s,\chi)\) at a first-order pole; \S 3.5 for the polar
divisor analysis.
\item[Arthur, J.] \emph{The Endoscopic Classification of
Representations}, AMS Colloquium~61, 2013. Chapter~1 \S 1.4--1.5
for the Arthur parameter of a residual representation and the
\(\mathrm{SL}_{2}\)-factor specification.
\item[Mo{\oe}glin, C.\ and Waldspurger, J.-L.] ``Le spectre
r{\'e}siduel de \(GL_{n}\)'', Ann.~Sci.~ENS~22 (1989), 605--674.
Arthur classification of residual representations for \(GL_{n}\):
pre-Arthur-2013 theorem restricted to \(GL_{n}\).
\item[Emerton, M.\ and Gee, T.] ``A geometric perspective on the
Breuil--M{\'e}zard conjecture'', arXiv:1908.07185 (2019). \S 6.2
unramified-locus smoothness; \S 6.3 closed point at the cyclotomic
character; the formal completion and tangent-space identification.
\item[Breuil, C.\ and M{\'e}zard, A.] ``Multiplicit{\'e}s modulaires
et repr{\'e}sentations de \(GL_{2}(\mathbb{Z}_{p})\) et de
\(\mathrm{Gal}(\overline{\mathbb{Q}}_{p}/\mathbb{Q}_{p})\) en
\(\ell=p\)'', Duke Math.~J.~115 (2002). Local deformation
multiplicities for two-dimensional residual Galois representations.
\item[Bella{\"i}che, J.\ and Chenevier, G.] \emph{Families of Galois
representations and Selmer groups}, Ast{\'e}risque~324, SMF 2009.
Chapter~2 for the residual gerbe at a reducible closed point.
\item[Schlessinger, M.] ``Functors of Artin rings'', Trans.~AMS~130
(1968). Pro-representability of the deformation functor.
\item[Mazur, B.] ``Deforming Galois representations'', in
\emph{Galois Groups over \(\mathbb{Q}\)}, MSRI Publ.~16 (1989).
Tangent space of the universal deformation ring.
\item[Bernstein, J.\ and Zelevinsky, A.] ``Induced representations
of reductive \(\mathfrak{p}\)-adic groups~I'', Ann.~Sci.~ENS~10
(1977). The composition factors of the principal series at
\(s=0,1\) and \(\mathrm{Ext}^{1}\)-groups between them at a
non-archimedean place.
\item[Wallach, N.] \emph{Real Reductive Groups~I}, Academic Press
1988. Chapter~5.4 for archimedean
\(\mathrm{Ext}^{1}(\mathbf{1},\mathrm{St})\) at \(GL_{2}(\mathbb{R})\).
\item[Jacquet, H.\ and Langlands, R.] \emph{Automorphic forms on
\(GL_{2}\)}, Lecture Notes in Math.~114, Springer 1970. \S 5 for the
non-splitting of the \(s=0,1\) induction at archimedean places.
\item[Harish-Chandra] ``Harmonic analysis on real reductive
groups~II'', Invent.~Math.~36 (1976). Theorem~1 for the archimedean
Plancherel density and residue.
\item[Kisin, M.] ``Potentially semi-stable deformation rings'',
J.~AMS~21 (2008). Cotangent complex of local Galois deformation
rings.
\item[Tate, J.] ``Fourier analysis in number fields and Hecke's
zeta functions'', Ph.D.~thesis, Princeton 1950; reprinted
Cassels--Fr{\"o}hlich 1967 Ch.~XV. Adelic restricted tensor
product; local-global factorisation.
\item[Positselski, L.] ``Two kinds of derived categories, Koszul
duality, and comodule-contramodule correspondence'',
arXiv:0905.2621 (2009). \S 6.3 extension of the bar--cobar pairing
to \(\mathrm{Ext}^{1}\)-classes.
\end{description}

\paragraph{Disjoint-source discipline.}
Chapter~\ref{v4-arith:w6-d:chapter} derivation uses
Mo{\oe}glin--Waldspurger~IV, Shahidi 2010 \S 3.5, Arthur 2013 \S 3.5;
the verification uses Mo{\oe}glin--Waldspurger~II.2,
Shahidi 2010 Thm~3.4.1, Arthur 2013 \S 1.4--1.5, plus
Emerton--Gee, Breuil--Mezard, Bella{\"i}che--Chenevier,
Schlessinger, Mazur, Kisin. The three shared authors
(Mo{\oe}glin--Waldspurger, Shahidi, Arthur) contribute through
disjoint subchapters on the two sides; all other sources are
unique to one side. primary-source independence holds.
